\section*{Abstract}
\begingroup\small
Does a direction that separates human and machine text also control the style of new generations? We test this question in paired base and instruction Qwen2.5-1.5B checkpoints using source-matched continuations, signed residual additions, projection ablation, five equal-norm random controls, and nuisance residualization. We evaluate 1,472 generated task outputs with a frozen encoder detector and a blinded external judge, alongside coherence, factual coverage, and unsupported claims. At the validation-selected layer, negative steering reduces judge-rated AI-likeness by 20.04 points on a 100-point scale (95\% paired bootstrap interval: $[-22.85,-17.27]$). However, coherence falls by 0.591 points on a five-point scale, and neither prespecified preservation criterion passes. The encoder does not corroborate the primary effect; positive steering increases judge-rated AI-likeness while decreasing its detector score. Readout is strong in both base and instruction backbones (held-out AUROC 0.9707 and 0.9704), but the instruction-backbone direction fails to transfer to base-generator continuations (0.486). Removing five nuisance proxies retains a smaller judged style effect with continued quality loss. These results establish partial control of perceived writing style under this protocol, while leaving a distinct, detector-consistent, content-preserving direction unestablished. They show why readout, controllability, measurement validity, and preservation require separate tests.
\par\endgroup
